% =====================================================================
% Section 11: MCQ feasibility, response-data availability, and inference cost
%
% ACM-format paper fragment. Intended to be \input under the MCQ section
% (the three \subsection blocks below sit under the paper's MCQ \section{}).
% A grouping \section is left commented out; uncomment if this file is
% included at top level.
%
% Authorship convention:
%   - Black prose is lightly cleaned from the outline PDF
%     "CAT for Accelerated Pedagogy Benchmarking.pdf"
%     (MCQ Feasibility; Availability and Quality of Response Data; Inference Cost).
%   - AI additions (numbers, analysis, figures/tables drawn from the repo)
%     are wrapped in \aiadd{...} and render blue.
%   - [CITATION NEEDED] markers render red.
%
% Requires (from the main preamble): graphicx, xcolor.
%
% Figure source paths (absolute), copied into ./figures/:
%   figures/diag_pedagogy_se0.3.png
%     <- /Users/arhant/Documents/EDLM/olmo-eval-full/AdaptiveTesting/Research/01_MCQ_ATLAS/data/pedagogy_feasibility/diag_pedagogy_se0.3.png
%   figures/calib_pedagogy_1pl_se0.3.png
%     <- /Users/arhant/Documents/EDLM/olmo-eval-full/AdaptiveTesting/Research/01_MCQ_ATLAS/data/pedagogy_feasibility/linear_calibration/calib_pedagogy_1pl_se0.3.png
%   figures/models_200_param_dist.png
%     <- /Users/arhant/Documents/EDLM/olmo-eval-full/AdaptiveTesting/Research/05_Data_Availability/figures/models_200_param_dist.png
%   figures/individual_latency_by_param_range.png
%     <- /Users/arhant/Documents/EDLM/olmo-eval-full/AdaptiveTesting/Research/04_Inference_Cost/data/individual_latency_by_param_range.png
%
% Table source data (absolute):
%   Table (MCQ CAT diagnostic)  <- .../pedagogy_feasibility/mcq_diagnostic_summary.csv
%                                  + .../pedagogy_feasibility/rasch_1pl/pedagogy_1pl_vs_2pl_summary.csv
%   Table (linear recalibration) <- .../pedagogy_feasibility/linear_calibration/pedagogy_linear_calibration.csv
%   Table (inference cost)      <- outline PDF; supported by .../04_Inference_Cost/data/individual_latency_by_param_range.csv
%
% Bibliography: refs/mcq_feasibility.bib
% =====================================================================

% Fallback macros: no-ops if the main preamble already defines them.
\providecommand{\aiadd}[1]{\textcolor{blue}{#1}}
\providecommand{\aibullets}[1]{{\color{blue}#1}}

% \section{ATLAS recreation: feasibility, data, and cost}

\subsection{Feasibility}

ATLAS shows that adaptive testing can cut the number of items needed to
predict benchmark accuracy \cite{li2025atlas}. There is still a gap between
that result and using evaluations during model training. Applying IRT
calibration to a benchmark with no public response data needs an honest
account of the upfront cost required to reach a target correlation. ATLAS
also did not show whether a benchmark can separate models at parameter
ranges where few calibration models are available.

We ran two feasibility experiments. The first calibrates the Pedagogy
benchmark, a skill-specific benchmark with no publicly available response
data \cite{lelievre2025pedagogy}. The second calibrates on subsets of ATLAS
data to imitate real-world conditions and to show the trade-offs between
model diversity, model count, and parameter-range match on calibration
correlation. \aiadd{The diversity, count, and range results are reported in
the experiments that follow. This subsection covers the pedagogy calibration
and two MCQ replications, using computerized adaptive testing built on Item
Response Theory \cite{lord1980irt,embretson2000irt}.}

\aiadd{We calibrated the Pedagogy CDPK set, 920 multiple-choice questions
\cite{lelievre2025pedagogy}, with the same bank calibration, EAP-theta
estimation, and Fisher-information item selection used in the ATLAS
replication. Fifty-two in-house models from the 156-model roster answered the
set; 40 were used for calibration and 12 were held out for validation (seed
7). After item screening, 318 of the 920 items carried usable response
variance. We report a 2PL bank and a 1PL/Rasch bank \cite{rasch1960} at two
stopping targets, standard error $\le 0.3$ and $\le 0.15$.}

\aiadd{Table~\ref{tab:mcq-cat} gives the result. At SE~$\le 0.3$ the 2PL bank
recovers the held-out ranking with Pearson $r = 0.72$ using about 10 of 318
items (3\%), but it over-predicts accuracy (MAE $0.063$). Tightening the
target to SE~$\le 0.15$ raises $r$ to $0.90$, but the test grows to about 276
items (87\%), which removes most of the saving. The 1PL/Rasch bank ranks
better at SE~$\le 0.3$ ($r = 0.89$) yet over-predicts the most (MAE $0.101$),
and at SE~$\le 0.15$ it never reaches the target, so it administers the full
bank.}

\aiadd{The two commonsense MCQ benchmarks behave very differently. PIQA
\cite{bisk2020piqa} replicates cleanly: at SE~$\le 0.3$ the 2PL bank reaches
$r = 0.94$ with about 11 of 831 items (1.3\%). SocialIQa
\cite{sap2019socialiqa} does not recover, with $r = 0.27$ at SE~$\le 0.3$ and
only $0.43$ at SE~$\le 0.15$, so adaptive selection cannot rank models on that
benchmark at either target. Feasibility depends on the benchmark and is not
automatic.}

\begin{table}[t]
\centering
\caption{\aiadd{MCQ CAT diagnostic on 12 held-out models (40 calibration
models, seed 7). Items kept after screening: pedagogy 318 of 920, PIQA 831 of
1{,}838, SocialIQa 969 of 1{,}954. ``Avg.\ items'' is the mean number
administered by the adaptive test; ``\% of bank'' is that as a share of the
kept items. Source: \texttt{mcq\_diagnostic\_summary.csv} and
\texttt{pedagogy\_1pl\_vs\_2pl\_summary.csv}.}}
\label{tab:mcq-cat}
\begin{tabular}{llcccrr}
\hline
Benchmark & Bank & SE~$\le$ & $r$ & MAE & Avg.\ items & \% of bank \\
\hline
Pedagogy  & 2PL & 0.30 & 0.716 & 0.063 & 10.2  & 3.2 \\
Pedagogy  & 2PL & 0.15 & 0.900 & 0.080 & 275.5 & 86.6 \\
Pedagogy  & 1PL & 0.30 & 0.890 & 0.101 & 43.8  & 13.8 \\
Pedagogy  & 1PL & 0.15 & 0.941 & 0.097 & 318.0 & 100.0 \\
PIQA      & 2PL & 0.30 & 0.937 & 0.040 & 10.6  & 1.3 \\
PIQA      & 2PL & 0.15 & 0.954 & 0.038 & 194.1 & 23.4 \\
SocialIQa & 2PL & 0.30 & 0.271 & 0.093 & 8.0   & 0.8 \\
SocialIQa & 2PL & 0.15 & 0.430 & 0.080 & 14.3  & 1.5 \\
\hline
\end{tabular}
\end{table}

\begin{figure}[t]
\centering
\includegraphics[width=\columnwidth]{figures/diag_pedagogy_se0.3.png}
\caption{\aiadd{Pedagogy CAT diagnostic, 2PL bank at SE~$\le 0.3$, on the 12
held-out models. Rank recovery is good ($r = 0.72$) with about 10 of 318
items, but every point sits above $y = x$: the raw prediction over-states
accuracy by a fixed offset. Source:
\texttt{pedagogy\_feasibility/diag\_pedagogy\_se0.3.png}.}}
\label{fig:diag-pedagogy}
\end{figure}

\aiadd{The pedagogy over-prediction is a fixed affine offset, not noise. The
IRT mean-probability prediction lives in a compressed range but is read
against the full 0 to 1 accuracy axis, while true pedagogy accuracy is
squeezed into about 0.22 to 0.31. The raw cloud has the right order but a
slope about six times too steep and a large positive intercept. A single
linear map fit on the 40 calibration models and applied to the 12 held-out
models removes almost all of the error and leaves $r$ unchanged
(Table~\ref{tab:pedagogy-recal} and Figure~\ref{fig:calib-pedagogy}). Honest
test-set MAE drops 76 to 90\% across the four cases. The 1PL/Rasch bank, the
worst raw case at MAE $0.101$, ends with the lowest calibrated MAE at about
$0.010$. Leave-one-out on the 12 test models agrees (0.008 to 0.016), so the
gain is real rather than fit noise.}

\begin{table}[t]
\centering
\caption{\aiadd{Pedagogy linear recalibration (honest train-fit: slope and
intercept fit on the 40 calibration models, applied to the 12 held-out
models). The map $\hat{y} = a\,x + b$ takes raw predicted accuracy $x$ to true
accuracy $\hat{y}$; $r$ is unchanged by an affine map. Source:
\texttt{pedagogy\_linear\_calibration.csv}.}}
\label{tab:pedagogy-recal}
\begin{tabular}{lccccrr}
\hline
Bank & SE~$\le$ & $r$ & MAE raw & MAE cal. & \% red. & $a$ / $b$ \\
\hline
2PL & 0.30 & 0.716 & 0.063 & 0.0151 & 76.1 & 0.152 / 0.214 \\
2PL & 0.15 & 0.900 & 0.080 & 0.0121 & 84.8 & 0.169 / 0.208 \\
1PL & 0.30 & 0.890 & 0.101 & 0.0104 & 89.7 & 0.171 / 0.208 \\
1PL & 0.15 & 0.941 & 0.097 & 0.0095 & 90.3 & 0.188 / 0.200 \\
\hline
\end{tabular}
\end{table}

\begin{figure*}[t]
\centering
\includegraphics[width=\textwidth]{figures/calib_pedagogy_1pl_se0.3.png}
\caption{\aiadd{Pedagogy 1PL/Rasch bank at SE~$\le 0.3$, before and after
linear recalibration on the 12 held-out models. Left: the raw prediction
tracks the ordering ($r = 0.89$) but rides a steep, shifted line (fit
$\mathrm{act} = 0.17\,\mathrm{pred} + 0.21$), giving MAE $0.101$. Right: after the
train-fit map the points fall onto $y = x$ at MAE $0.010$, a 90\% reduction,
with $r$ unchanged. Source:
\texttt{linear\_calibration/calib\_pedagogy\_1pl\_se0.3.png}.}}
\label{fig:calib-pedagogy}
\end{figure*}

\aiadd{For pedagogy the practical recipe is a 1PL/Rasch bank at SE~$\le 0.3$
for ranking, plus a linear recalibration step when an absolute accuracy number
is needed. This recovers the ranking with roughly 3 to 14\% of the bank,
depending on the bank, and brings predicted accuracy within about one point of
the full-bank value.}

\subsection{Availability and quality of response data}

Two of the higher-quality sources of per-question model responses are the Open
LLM Leaderboard (OpenLM) \cite{beeching2023openllm,fourrier2024openllm} and
HELM (OpenHELM) \cite{liang2023helm}. ATLAS also uses OpenLM. Prompting
strategy and scoring vary across these sources, and new benchmarks take time
to be added to OpenLM.

\aiadd{We checked five sources for reusable response data. Four are available
to us; OpenHELM is not present in our local copies and is referenced here only
for context. The Open LLM Leaderboard v2 dumps \cite{fourrier2024openllm}
cover BBH, GPQA, MATH, IFEval, and MuSR for 1{,}102 models spanning 0.213 to
7.0B parameters (about 28~GB), keeping the question, response text,
prediction, gold answer, result, and scoring method for each item. Of these,
1{,}101 downloaded cleanly, 84 had no sample files, and 3 errored. The ATLAS
release \cite{li2025atlas} provides a full 0/1 response matrix only for ARC
\cite{clark2018arc} (3{,}747 calibration and 417 held-out models); its other
benchmarks ship item parameters and accuracy summaries rather than full
matrices.}

\aiadd{Quality and coverage vary in ways that matter for pooling. OpenLM mixes
scoring methods (\texttt{acc\_norm} for BBH, GPQA, and MuSR; \texttt{exact\_match}
for MATH; a strict prompt score for IFEval) and keeps the raw generated text,
while our in-house MCQ dumps reduce each item to a loglikelihood choice with no
latency. Any cross-source pooling has to reconcile these first. The two large
sources also barely overlap: OpenLM carries the newer benchmarks but none of
the five ATLAS core MCQ sets (ARC, HellaSwag, GSM8K, TruthfulQA, WinoGrande),
and a skill-specific benchmark such as Pedagogy has no public response data at
all. That absence is why the feasibility study above had to generate its own
response matrix.}

\aiadd{The curated roster (Figure~\ref{fig:paramdist}) holds 156 checkpoints
from 0.218 to 7.0B parameters (mean 3.6B, median 3.0B) across 50
organizations. The sizes cluster at 1 to 2B (42 models) and at 6 to 7B (52
models), the 5 to 6B band is empty, and the 2 to 5B middle is thin. This shape
matters for calibration: the ATLAS recalibration pool skews even harder, with
about 95\% of its 1{,}686 models at 7B, so mid-range coverage is the weak
point when matching a calibration set to a target parameter range.}

\begin{figure*}[t]
\centering
\includegraphics[width=\textwidth]{figures/models_200_param_dist.png}
\caption{\aiadd{Parameter distribution of the 156-model calibration roster
(left) and the top organizations by model count (right). Coverage clusters at
1 to 2B and at 6 to 7B, with a thin 2 to 5B middle and no 5 to 6B models.
Source: \texttt{05\_Data\_Availability/figures/models\_200\_param\_dist.png}.}}
\label{fig:paramdist}
\end{figure*}

\aiadd{New benchmarks are also slow to appear in OpenLM. Its coverage is fixed
to the v2 task set, so a skill-specific benchmark such as Pedagogy is absent
and has to be run in-house. For a team that depends on public response data,
this means the calibration data for a new skill does not exist yet, which
pushes the work onto local inference and makes the inference cost below the
deciding factor.}

\subsection{Inference cost}

The inference cost is reported for the 0 to 7B parameter range. L4 GPUs were
used for inference, and CPUs parallelized the model downloads into an
accessible cache.

\begin{table}[t]
\centering
\caption{Inference cost for tutor models, reproduced from the outline.
\aiadd{The header label is corrected below: the measurements are an L4
estimate, not L40S.}}
\label{tab:inference-cost}
\begin{tabular}{cc}
\hline
Parameter range & Average individual inference \\
(billions of parameters) & time on L40S GPU (seconds) \\
\hline
0-1 & 4 \\
1-2 & 10 \\
2-5 & 19 \\
5-7 & 34 \\
\hline
\end{tabular}
\end{table}

\aiadd{Two things about Table~\ref{tab:inference-cost} need stating plainly.
First, the numbers are single-stream latency, the time to answer one request
at a time, and they are an L4 memory-bandwidth roofline estimate (300~GB/s,
60\% efficiency, 2 bytes per parameter), not a direct measurement. Second, the
header says L40S but the inference actually ran on an L4, a g6.xlarge instance
at \$0.8048 per hour. An L40S has roughly 864~GB/s of bandwidth and would
decode about two to three times faster, so the L40S label understates the cost
on the hardware that was used. The table should be read as L4.}

\aiadd{In production the roster ran batched under vLLM \cite{kwon2023vllm},
where the measured per-request time is 0.25 to 1.6 seconds across the same
size bands, a 14 to 21 times speed-up over the single-stream estimate
(Figure~\ref{fig:latency}). Which number applies depends on the question: a
cold, one-request-at-a-time diagnostic sees the 4 to 34 second figures, while a
warm batched server at checkpoint time sees the sub-two-second figures.
Converted to dollars, an MCQ inference costs about \$0.009 per 1{,}000
requests, and an open-ended inference costs about \$0.20 to \$0.23 per 1{,}000,
roughly 20 to 25 times more, because open-ended cost is driven by output
tokens rather than input length.}

\begin{figure*}[t]
\centering
\includegraphics[width=0.75\textwidth]{figures/individual_latency_by_param_range.png}
\caption{\aiadd{Single-stream (roofline estimate) versus batched vLLM
(measured) latency per open-ended inference on an L4 GPU, by parameter range.
Batching cuts per-request time 14 to 21 times. The single-stream bars are the
L4 roofline estimates (about 4, 10, 20, and 35 seconds) that underlie the
outline's tutor-cost table.
Source: \texttt{04\_Inference\_Cost/data/individual\_latency\_by\_param\_range.png}.}}
\label{fig:latency}
\end{figure*}

There are two points of discussion. The first is the reduction in time during
inference, which allows CAT to be used at checkpoints and speeds up
experimentation. The second is the actual delta in cost.

\aiadd{On the first point, the adaptive test needs only a small fraction of the
bank once items are calibrated: about 3\% of the pedagogy items with the 2PL
bank and about 14\% with the 1PL bank at SE~$\le 0.3$, and about 1\% of PIQA.
With batched inference at well under two seconds per request, the marginal cost
of running a diagnostic on a fresh checkpoint is small enough to run during
training rather than only at the end.}

\aiadd{On the second point, the cost is front-loaded. The large expense is the
one-time calibration pass, running the full model roster over the full item
bank, plus the judge grading for open-ended tasks. After that, each new model
pays only for the items its adaptive test actually selects, roughly 7 to 30
times fewer than a full pedagogy pass, depending on the bank. The upfront cost
buys cheap, repeatable evaluation afterward, which is what makes
checkpoint-time testing practical.}
